\section{壁②：時系列は本当にAIレディか}

% =====================================================================
%  Frame 1 : 一見、時系列はAIレディ
% =====================================================================
\begin{frame}{一見、時系列は「もうAIレディ」に見える}{株価・金利・出来高… アクセスも分析もできる}
  \begin{columns}[T]
    \begin{column}{0.5\textwidth}
      \begin{itemize}\small
        \item 株価・金利・為替・出来高――数値の時系列は、いま最も\kw{整備が進んだ}データに見える。
        \item だから素朴には、こう思ってしまう。
      \end{itemize}
      \vskip2mm
      \begin{casestudy}{よくある見方}
        \small
        「アクセスもできる、APIもある、分析もできる。\\
        時系列は\kw{もう解決済み}なのでは？」
      \end{casestudy}
    \end{column}
    \begin{column}{0.46\textwidth}
      \begin{center}
      \scalebox{0.72}{%
      \begin{tikzpicture}[node distance=2.4mm,font=\footnotesize]
        \node[fnode good,text width=5.4cm] (a)
          {\textbf{✓ アクセスできる}\\ヒストリカルデータが揃う};
        \node[fnode good,below=of a,text width=5.4cm] (b)
          {\textbf{✓ APIがある}\\J-Quants等で機械的に取得};
        \node[fnode good,below=of b,text width=5.4cm] (c)
          {\textbf{✓ 分析もできる}\\可視化・統計量・バックテスト};
        \node[fnode good,below=of c,text width=5.4cm] (d)
          {\textbf{✓ LLMも読める}\\数列をトークンとして入力可};
        \draw[farr] (a)--(b);
        \draw[farr] (b)--(c);
        \draw[farr] (c)--(d);
      \end{tikzpicture}}
      \end{center}
      \figcap{「できていること」は多い。だが十分だろうか。}
    \end{column}
  \end{columns}
\end{frame}

% =====================================================================
%  Frame 2 : 【落とし穴】本当に"理解"しているのか
% =====================================================================
\begin{frame}{【落とし穴】本当に「理解」しているのか}{読めること $\neq$ 統計的構造を内在化していること}
  \begin{columns}[T]
    \begin{column}{0.46\textwidth}
      \begin{itemize}\footnotesize\setlength\itemsep{1.0mm}
        \item LLMは時系列を\kwbad{トークン列としては読める}が、\\
          非定常性・レジーム変化・\kwbad{データリーク（先読み）}・因果を内在化していない。
        \item 学習カットオフ\emph{以降}に評価をずらすだけで――
      \end{itemize}
      \vskip0.4mm
      \begin{insight}{ベンチのアルファは脆い}
        \footnotesize
        77件のLLMトレーディング研究のレビューで、取引コストを明示モデル化したのは\kwbad{わずか1件}、時間整合的で再現可能な評価は\kwbad{わずか2件}のみ。
      \end{insight}
    \end{column}
    \begin{column}{0.5\textwidth}
      \begin{center}
      \scalebox{0.9}{%
      \begin{tikzpicture}[font=\footnotesize,y=0.12cm]
        % 縦軸目盛（0 と 20）とゼロ基準線
        \draw[brandline,line width=0.4pt] (0,0) -- (6.2,0);
        \draw[brandline,line width=0.4pt] (0,20) -- (6.2,20);
        \node[left,text=brandmute,font=\scriptsize] at (0,0) {0};
        \node[left,text=brandmute,font=\scriptsize] at (0,20) {20};
        \draw[brandmute,line width=0.8pt] (0,-4) -- (0,24);
        % バー：カットオフ前（緑・+20.73）
        \fill[brandgood!75,draw=brandgood!55!black,line width=0.6pt]
          (0.9,0) rectangle (2.3,20.73);
        \node[above,text=brandgood!45!black,font=\footnotesize\bfseries] at (1.6,20.73) {+20.73\%};
        \node[below,text=brandgray,font=\scriptsize] at (1.6,0) {カットオフ前};
        % バー：カットオフ後（赤・-1.04）
        \fill[brandbad!78,draw=brandbad!55!black,line width=0.6pt]
          (3.9,0) rectangle (5.3,-1.04);
        \node[below=1mm,text=brandbad,font=\footnotesize\bfseries] at (4.6,-1.04) {-1.04\%};
        \node[above,text=brandgray,font=\scriptsize] at (4.6,0.4) {カットオフ後};
        \node[rotate=90,text=brandmute,font=\scriptsize] at (-0.75,10) {年率アルファ (\%)};
        \draw[farr accent] (2.4,18) to[bend left=18] (3.8,2);
      \end{tikzpicture}}
      \end{center}
      \figcap{カットオフ以降に評価をずらすと年率アルファ \kwbad{+20.73\% → -1.04\%}。}
      \srcnote{Look-Ahead-Bench, 2026 / Agentic trading review, 2026}
    \end{column}
  \end{columns}
\end{frame}

% =====================================================================
%  Frame 3 : なぜ難しいか
% =====================================================================
\begin{frame}{なぜ難しいか：意味は「数列」に宿らない}{時系列の意味は、パラメータ・モデル・因果構造に宿る}
  \begin{columns}[T]
    \begin{column}{0.48\textwidth}
      \begin{itemize}\small
        \item 時系列の意味は「数値の羅列」そのものではなく、\\
          その背後の\kw{パラメータ・モデル・因果構造}に宿る。
        \item LLMはテキストとしては処理できても、\\
          \kwbad{統計的構造を理解しているわけではない}。
      \end{itemize}
      \vskip1.5mm
      \begin{keytake}{壁①と同じ罠が、より深い形で}
        \small
        壁①の教訓「\kw{読める $\neq$ 使える}」が、\\
        時系列では\kwg{さらに顕著}に立ちはだかる。
      \end{keytake}
    \end{column}
    \begin{column}{0.48\textwidth}
      \begin{center}
      \scalebox{0.85}{%
      \begin{tikzpicture}[font=\footnotesize]
        \node[fnode muted,text width=3.0cm] (raw)
          {\textbf{数列そのもの}\\ \texttt{102, 98, 105,\,…}};
        \node[fnode blue,right=1.4cm of raw,text width=3.2cm] (mean)
          {LLMが見るもの\\= \textbf{トークン列}};
        \node[fnode accent,below=1.3cm of raw,text width=3.0cm] (param)
          {\textbf{パラメータ}\\ボラ・自己相関・レジーム};
        \node[fnode accent,below=1.3cm of mean,text width=3.2cm] (causal)
          {\textbf{因果構造}\\何が何を動かすか};
        \draw[farr dash] (raw)--(mean);
        \node[text=brandbad,font=\scriptsize\bfseries] at ($(raw)!0.5!(mean)+(0,0.5)$) {表層のみ};
        \draw[farr accent] (raw)--(param);
        \draw[farr accent] (mean)--(causal);
        \node[text=brandaccent,font=\scriptsize\bfseries,align=center]
          at ($(param)!0.5!(causal)+(0,-1.0)$) {意味はここに宿る};
      \end{tikzpicture}}
      \end{center}
      \figcap{LLMは表層のトークン列に届くが、意味の在処には届かない。}
    \end{column}
  \end{columns}
\end{frame}

% =====================================================================
%  Frame 4 : 【中心図】解＝専門AIエージェント＋モデルのブリッジ
% =====================================================================
\begin{frame}{解：専門AIエージェントがモデルを解釈し、LLMへブリッジする}{時系列 → 専門エージェント → 解釈 → ブリッジ → アウトカム}
  \begin{center}
  \scalebox{0.75}{%
  \begin{tikzpicture}[font=\footnotesize,node distance=6mm]
    \node[fnode,text width=2.5cm] (ts)
      {\textbf{時系列データ}\\株価・金利・出来高};
    \node[fnode accent,right=9mm of ts,text width=3.0cm] (agent)
      {\textbf{専門AIエージェント}\\統計・DL・時系列モデルを\\\emph{道具として実行}};
    \node[fnode blue,right=9mm of agent,text width=2.9cm] (interp)
      {\textbf{パラメータ・\\意味の解釈}\\状態空間・GARCH・因子};
    \node[fnode blue,right=9mm of interp,text width=2.7cm] (bridge)
      {\textbf{LLMが読める形\\へブリッジ}\\記号・言語化};
    \node[fnode dark,right=9mm of bridge,text width=2.3cm] (out)
      {\textbf{アウトカム}\\（リターン）};
    \draw[farr] (ts)--(agent);
    \draw[farr accent] (agent)--(interp);
    \draw[farr blue] (interp)--(bridge);
    \draw[farr blue] (bridge)--(out);
    % 専門エージェントが使う道具（下に添える）
    \node[fnode muted,below=10mm of interp,text width=8.6cm,font=\scriptsize] (tools)
      {\textbf{道具箱：}状態空間モデル／GARCH／因子モデル／ディープラーニング（時系列）――\;
       エージェントが選び、当てはめ、\kw{パラメータの意味}を取り出す};
    \draw[farr dash] (agent.south) to[bend right=12] (tools.west);
    \draw[farr dash] (tools.east) to[bend right=12] (interp.south);
  \end{tikzpicture}}
  \end{center}
  \begin{center}
    \bigstate{「数列」ではなく\kwg{解釈済みの構造}を、LLMに手渡す}
  \end{center}
\end{frame}

% =====================================================================
%  Frame 5 : それでもまだAIレディでない
% =====================================================================
\begin{frame}{それでも、まだAIレディではない}{フレームワーク（部品）はあるが、一個一個つながっていない}
  \begin{columns}[T]
    \begin{column}{0.5\textwidth}
      \begin{itemize}\wide
        \item モデルも、エージェントも、検証手法も\\
          ――\kw{部品（フレームワーク）は揃っている}。
        \item しかし、それらが\kwbad{一個一個つながっていない}。\\
          これが壁①より\kwg{深い落とし穴}。
        \item 解決には、意思決定まで\kw{一気通貫}でつなぎ、\\
          \kwg{検証層}を組み込むことが要る。
      \end{itemize}
    \end{column}
    \begin{column}{0.48\textwidth}
      \begin{insight}{実装条件：検証層を組み込む}
        \footnotesize
        \begin{itemize}\setlength\itemsep{1mm}
          \item \kwg{棄権}――「答えてはいけない状況」を検知する能力。
          \item グラフ・記号での\kw{経路検証}（FinChain / GBFR）。
          \item \kw{反事実・因果}を組み込んだ検証（FinCARE ほか）。
        \end{itemize}
        \vskip1mm
        {\color{brandgray} 金融AIの信頼性は、当てる力より\kwg{「答えない力」}に宿る。}
      \end{insight}
    \end{column}
  \end{columns}
\end{frame}

% =====================================================================
%  Frame 6 : 壁②まとめ
% =====================================================================
\begin{frame}{壁②まとめ：時系列は「理解」でつまずく}
  \begin{keytake}{壁②の要点}
    \begin{itemize}\setlength\itemsep{2mm}
      \item 時系列は\kw{「アクセスできる」}のと\kwbad{「理解している」}のがまったく別物。\\
        {\small 学習カットオフをまたぐだけで、アルファは \kwbad{+20.73\% → -1.04\%} へ消える。}
      \item 解は、\kwg{専門AIエージェント}が統計・DL・時系列モデルを道具として実行し、\\
        その\kw{パラメータと意味を解釈}して、\kw{LLMが読める形にブリッジ}すること。
      \item ただし部品はまだ\kwbad{つながっていない}。\\
        \kwg{棄権・反事実・因果}を組み込んだ検証層こそ、実装の条件。
    \end{itemize}
  \end{keytake}
  \vskip1mm
  \begin{center}
    \bigstate{時系列のAIレディ化とは、\kwg{解釈と検証を一気通貫でつなぐ}こと}
  \end{center}
\end{frame}
